% Source-derived chapter 03: Models for the arc space
% Original source: intersection-complexes-unramified-l-factors
\chapter{Models for the arc space}
\label{intersection-complexes-unramified-l-factors-chapter-03}
\source{intersection-complexes-unramified-l-factors}

\begin{remark}[Source section]
\label{intersection-complexes-unramified-l-factors-chapter-03-source}
\source{intersection-complexes-unramified-l-factors}
\source{intersection-complexes-unramified-l-factors}
This chapter reproduces the corresponding section of the retrieved
LaTeX source, including its definitions, statements, examples, and
proofs. It has not yet been translated into Lean.
\notready
\end{remark}

% The source section title was: Models for the arc space
\label{sect:models}
\source{intersection-complexes-unramified-l-factors}

In this section we define two models (in the sense of Grinberg--Kazhdan) 
for the arc space of $X$, both of which
were already introduced in \cite{GN} and go back to ideas of Drinfeld. 
We call these the global and Zastava models (the term `global' refers to the fact
that the model depends on the curve $C$). 
The global model $\sM_X$ is crucial because it allows us to model the $G(\mf o)$-orbits of $X(\mf o)$,
something which cannot be done directly via the Zastava model. 
On the other hand, the Zastava model $\Scr Y_X$ is more suitable for finite type calculations. 

\smallskip

Various incarnations of these constructions have been used in 
\cite{FM, FFKM, BG, BFGM, BFG, ABB}.
To place our work in this context, we remark that when $X = \ol{N\bs G}^\aff$ 
we have $\sM_X = \ol\Bun_N$ is Drinfeld's compactification of $\Bun_N$ 
and $\sY_X$ is ``the'' Zastava space of \cite{FM}. % (in this case $\ch G_X = \ch T$).

\smallskip

While Gaitsgory--Nadler define the global and Zastava models for any affine $X$, 
in order to avoid various technical difficulties they faced (such as the existence
of \emph{twisted strata}, which are related to the existence of 
disconnected stabilizer subgroups) we make the following simplifying assumption:

\medskip

Starting from \S\ref{sect:localmodel}, 
\emph{we assume for the rest of the paper that $B$ acts simply transitively on $X^\circ$.} 
If $\ch G_X = \ch G$ and $X$ has no spherical roots of type $N$, then the assumption above always holds. 


\begin{rem}
\source{intersection-complexes-unramified-l-factors}
\notready \label{rem:Hconnected}
\source{intersection-complexes-unramified-l-factors}
The assumption that $B$ acts simply transitively on $X^\circ$ implies that $H$ must
be connected, by Lemma~\ref{lem:H_I-conn}. 
\end{rem}


\subsection{Global model}

Gaitsgory--Nadler \cite{GN} define certain stacks of meromorphic quasimaps from $C \dashrightarrow X/G$ 
to model $X^\bullet(F)$, the loop space of $X^\bullet$. 
Our global model $\sM_X$ is the substack consisting of those quasimaps that 
extend to regular maps\footnote{In the literature, when $X=\ol{N\bs G}$ regular maps 
$C \to \ol{N\bs G}$ are still referred to as quasimaps to $N\bs G$.} $C \to X/G$. 

\subsubsection{Definition}
Define the stack 
\[ \sM_{X} := \Maps_\gen(C, X/ G \supset X^\bullet /G) \]
to be the open substack of $\Maps(C, X/G)$ representing maps generically
landing in $X^\bullet/G = (H\bs G)/G = H\bs\pt$. 

\smallskip

An $S$-point of $\Scr M_{X}$ is a map $f : C\times S \to X/G$, which is
equivalent to the datum $(\Scr P_G, \sigma)$ where 
\begin{itemize}
\item $\Scr P_G$ is a $G$-bundle on $C\times S$ and 
\item $\sigma : C\times S \to X\times^G \Scr P_G$ is a section over $C\times S$ 
such that
\item $\sigma|_{\Spec k(C) \xt S}$ lands in $X^\bullet \xt^G \Scr P_G = H\bs \Scr P_G$
and gives $\Scr P_G |_{\Spec k(C)\xt S}$ the structure of an $H$-bundle.
\end{itemize}
We call the preimage $\sigma^{-1}(X^\bullet \xt^G \sP_G)\subset C \xt S$ the 
locus of $G$-nondegenerate points. 

\begin{prop}
\source{intersection-complexes-unramified-l-factors}
\notready
The natural map $\Scr M_{X} \to \Bun_G$ is schematic locally of finite type. 
In particular, $\Scr M_{X}$ is an algebraic stack locally of finite type over $k$.
\end{prop}
\begin{proof}
Since $\Scr M_{X}$ is an open substack of $\Maps(C, X/G)$,
it suffices to show the latter is schematic locally of finite type
over $\Bun_G$. 
Let $S \to \Bun_G$ correspond to a $G$-bundle $\Scr P_G$ on $C\xt S$.
Then the fiber product $S \underset{\Bun_G}\times \Maps(C, X/G)$ 
is isomorphic to the space of sections $C\xt S \to X \times^G \Scr P_G$ over $C \xt S$.
This space is representable by a $k$-scheme locally of finite type 
(\cite[Theorem 5.23]{FGA-explained}). 
\end{proof}

If $X = H\bs G$ is homogeneous, then $\Scr M_X = \Maps(C,H\bs \pt) = \Bun_H$. 

\subsubsection{Adelic description} \label{sect:adelic}
\source{intersection-complexes-unramified-l-factors}
For motivational purposes, we give an ``adelic'' description of the $k$-points of 
$\sM_X$. 
Let $\mbb A$ denote the restricted product $\prod'_{v\in \abs C} F_v$ and 
$\mbb O = \prod_{v\in \abs C} \mf o_v$. 
(If $k=\mbb F_q$ then $\mbb A$ is the ring of adeles of the function field $\Bbbk = k(C)$.)

The underlying set of meromorphic quasimaps $C \dashrightarrow X/G$ 
can be identified with the set
\[ X^\bullet(\Bbbk) \xt^{G(\Bbbk)} G(\mbb A) / G(\mbb O) 
= H(\Bbbk) \bs G(\mbb A) / G(\mbb O). 
\]
Note that the $G$-action on $X$ induces a map 
\[ X^\bullet(\Bbbk) \xt^{G(\Bbbk)} G(\mbb A) / G(\mbb O) \to X^\bullet(\mbb A)/G(\mbb O) \subset X(\mbb A)/G(\mbb O). 
\]
The underlying set of $\sM_X(k)$ identifies with 
the preimage of $(X^\bullet(\mbb A) \cap X(\mbb O)) / G(\mbb O)$ under the map above.

\medskip

Note that the topologies on $X^\bullet(\mbb A)$ vs.~$X(\mbb A)$ are different: the
fact that the geometric constructions above depend on $X$ can be expressed
by saying that we are always using the topology of $X(\mbb A)$, not of $X^\bullet(\mbb A)$.



\subsubsection{} \label{def:partition}
\source{intersection-complexes-unramified-l-factors}
For any set $\msf S$, define the set of unordered \emph{multisets} in
$\msf S$ to be the formal direct sum
\[ \Sym^\infty(\msf S) := \bigoplus_{\ch\lambda\in \msf S} \mbb N[\ch\lambda]. \] 
An element of $\Sym^\infty(\msf S)$ is a 
formal sum $\mf P = \sum_{\ch\lambda \in \msf S} N_{\ch\lambda}[\ch\lambda]$
where $N_{\ch\lambda}\ge 0$ are integers and only finitely many 
are nonzero. 
For a multiset $\mf P$, define 
\[ \oo C^{\mf P} = \oo \prod_{\ch\lambda \in \msf S} \oo C^{(N_{\ch\lambda})}
\]
to be the open subscheme of $C^{\mf P}:=\prod C^{(N_{\ch\lambda})}$ 
with all diagonals
removed, i.e., the subscheme of multiplicity free divisors.
We write $\mf P=0$ for the zero element, and we will use the 
convention $\oo C^{0}=C^{0}=\pt$. 
If we define $\abs{\mf P} = \sum N_{\ch\lambda}$, then there are 
natural maps $C^{\abs{\mf P}} \to C^{\mf P} \to C^{(\abs{\mf P})} \subset \Sym C$.

Now consider the case when $\msf S = \msf M\sm 0$ for a commutative monoid
$\msf M$. Then $\Sym^\infty(\msf M\sm 0)$ identifies with the set of 
\emph{partitions} of arbitrary elements in $\msf M$.
Given a partition $\mf P\in \Sym^\infty(\msf M\sm 0)$ as above, 
we define $\deg : \Sym^\infty(\msf M\sm 0) \to \msf M$ by 
\[
    \deg(\mf P) := \sum_{\ch\lambda\in \msf M\sm 0} N_{\ch\lambda}\ch\lambda 
\]
with addition taking place in $\msf M$. 
There is a natural order on the set $\Sym^\infty(\msf M\sm 0)$: 
we say that $\mf P$ refines $\mf P'$ if the difference $\mf P-\mf P'$
viewed as an element of $\bigoplus_{\ch\lambda\in \msf M\sm 0}
 \mbb Z [\ch\lambda]$ can be written as a sum of elements of the
form $[\ch\lambda']+[\ch\lambda'']-[\ch\lambda]$ with $\ch\lambda'+\ch\lambda''=\ch\lambda$ in $\msf M$. 

Let $\Prim(\msf M)$ be the set of primitive elements of $\msf M$, i.e.,
the elements $\ch\lambda\in \msf M\sm 0$ that cannot be decomposed
as a sum $\ch\lambda=\ch\lambda_1+\ch\lambda_2$ where $\ch\lambda_1,\ch\lambda_2 \in \msf M\sm 0$. Then any partition in $\Sym^\infty(\msf M\sm 0)$
can be refined to an element in $\Sym^\infty(\Prim(\msf M))$. 



\subsubsection{Stratification of $\sM_X$}

\label{sect:globstrat}
\source{intersection-complexes-unramified-l-factors}


We would like to stratify $\sM_X$ according to $G(\mf o_v)$-orbits of 
$X(\mf o_v)\cap X^\bullet(F_v)$ at each $v\in \abs C$, described in Theorem \ref{thm:Gorbits}. 


Consider the set $\Sym^\infty(\mathfrak c_X^-\sm 0)$
of partitions in $\mathfrak c_X^-$, as defined
in \S\ref{def:partition}. 
Let  
\[ \ch\Theta = \sum_{\ch\theta\in \mathfrak c_X^- \sm 0} N_{\ch\theta} [\ch\theta] \] 
denote such a partition. 
We will write $\ch\Theta=0$ for the empty partition and 
$\ch\Theta=[\ch\theta]$ for the singleton partition corresponding to a
single element $\ch\theta$. 

In \S\ref{proof:globstrata}, we define locally closed substacks $\Scr M^{\ch\Theta}_X$ of $\Scr M_X$ ranging over all partitions 
$\ch\Theta$ in $\mathfrak c_X^-$.  
For simplicity we only describe $\Scr M^{\ch\Theta}_X$ on $k$-points below:
Such a point consists of a map $f : C\to X/G$ and a formal sum
$\sum_{v\in \abs C} \ch\theta_v \cdot v$ satisfying the following
conditions: 
\begin{itemize}
\item $\ch\theta_v \ne 0$ for finitely many $v\in \abs C$, 
\item for a fixed $\ch\theta \ne 0$, the cardinality of $\{ v\in \abs C \mid \ch\theta_v =\ch\theta\}$ equals $N_{\ch\theta}$, 
\item for each $v\in \abs C$
the restriction $f|_{\Spec \mf o_v} : \Spec \mf o_v \to X/G$ defines a point in $\msf L^{\ch\theta_v} X / \msf L^+ G$.  
\end{itemize} 

\begin{lem}
\source{intersection-complexes-unramified-l-factors}
\notready \label{lem:globstrata}
\source{intersection-complexes-unramified-l-factors}
The substack $\Scr M^{\ch\Theta}_X$ is
smooth and locally closed in $\Scr M_X$. 
\end{lem}

We defer the proof of Lemma~\ref{lem:globstrata} to \S\ref{proof:globstrata} of the appendix. 

\medskip
\begin{prop}
\source{intersection-complexes-unramified-l-factors}
\notready \label{prop:globwhitney}
\source{intersection-complexes-unramified-l-factors}
Assume $k=\mbb C$. Let $\Scr S$ denote the collection of connected components of 
$\Scr M^{\ch\Theta}_X$, ranging over all 
$\ch\Theta \in \Sym^\infty(\mathfrak c_X^- \sm 0)$. Then $\Scr S$ 
is a Whitney\footnote{We say that a stratification on an algebraic stack $\Scr M$ locally of
finite type is Whitney if the stratification is Whitney after pullback
along any (equivalently all) smooth cover of $\Scr M$ by a scheme.} 
stratification of $\Scr M_X$. 
\end{prop}

By a stratification we mean a collection of locally closed substacks that form a
disjoint union on $k$-points and such that the closure of any stratum is a union of
strata. The stratification is Whitney if the strata are smooth and
every pair of strata satisfies Whitney's condition B. (This only makes sense
if the characteristic of $k$ is zero; in positive characteristic, see 
\S \ref{sect:whitneyperverse} below.)

\smallskip

The proof of Proposition~\ref{prop:globwhitney} is given in \S\ref{proof:whitney}. 
We call this the \emph{fine stratification} of $\Scr M_X$. 
The open stratum $\sM_X^0 = \Maps(C, H\bs G/G)$ identifies
with $\Bun_H$. We abbreviate $\sM_X^{\ch\theta} := \sM_X^{[\ch\theta]}$. 

\begin{rem}
\source{intersection-complexes-unramified-l-factors}
\notready
If $X$ is affine and homogeneous, then $\Cal C_0(X) \cap \Cal V = 0$ 
so the stratification is trivial, consisting of
the single smooth stratum $\Scr M_X$ itself.
\end{rem}

\subsubsection{} \label{sect:whitneyperverse}
\source{intersection-complexes-unramified-l-factors}
Let $\rmD^b_{\Scr S}(\sM_X,\ol\bbQ_\ell)$  
denote the subcategory of bounded $\Scr S$-constructible complexes, i.e., 
the usual cohomology sheaf $\rmH^i(\Scr F)|_S$ is a local system
of finite rank for all $i\in\bbZ$ and $S\in\Scr S$. 
As explained in \cite{BBDG}, the category $\rmD^b_{\Scr S}(\sM_X)$ has a 
perverse $t$-structure. Let $\rmP_{\Scr S}(\sM_X)$ denote the heart
of this $t$-structure, i.e., this is the abelian subcategory of all perverse sheaves
that are $\Scr S$-constructible. 

In particular, the IC complex of the closure of any stratum $\sM_X^{\ch\Theta}$ 
is an object of $\rmP_{\Scr S}(\sM_X)$.
When $k$ has positive characteristic, this is the condition on $\Scr S$ 
that we need (in place of the Whitney condition). 
Proposition~\ref{prop:Whitney-poschar} shows that this condition 
indeed holds in positive characteristic.


\subsection{Toric case} \label{sect:baseA}
\source{intersection-complexes-unramified-l-factors}
If we apply the definitions above to the special case where $G$ is replaced by the torus $T_X$ and $X$ is replaced by the toric variety $X\sslash N$, we obtain the space \begin{equation}  \label{e:defA}
\source{intersection-complexes-unramified-l-factors}
    \Scr A = \Maps_\gen(C, (X/\!\!/N)/T_X \supset \pt) 
\end{equation}
of maps generically landing in $T_X/T_X=\pt$. 



The stack $\Scr A$ has been previously studied in \cite[\S 3]{BNS} as 
a model for the formal arc space of the toric variety $X/\!\!/N$ 
(in particular, $\Scr A$ turns out to be representable by a scheme). 
We review the relevant properties below.

For any $N \in \mbb N$, we have the $N$th symmetric product
$C^{(N)}$ of $C$, which identifies with the Hilbert scheme $\on{Hilb}^N(C)$
parametrizing relative effective divisors in $C$ of degree $N$.
Let $\Sym C$ denote the disjoint union 
$\bigsqcup_{N\in \mathbb N} C^{(N)}$ (where $C^{(0)}=\pt$). 


\begin{eg}
\source{intersection-complexes-unramified-l-factors}
\notready
Observe that $\Maps_\gen(C, \mbb A^1/\mbb G_m\supset \pt)$ sends a test scheme $S$ to 
the set of relative effective Cartier divisors on $C \times S$, i.e., 
$\Maps_\gen(C, \mbb A^1/\mbb G_m\supset \pt) \cong \Sym C$. 
Addition of divisors gives $\Sym C$ the structure of a monoid. 
\end{eg}


Let $\mf c_X^\vee = \Hom(\mf c_X, \mbb N)$ denote the monoid dual 
to $\mf c_X$, so $k[X/\!\!/N]$ is the semigroup algebra of $\mf c_X^\vee$. 
Then there is an isomorphism 
\[ \Scr A \cong \Hom(\mf c_X^\vee, \Sym C) \] 
where the right hand side represents homomorphisms of monoid objects in the category of schemes
(with $\mf c_X^\vee$ viewed as a discrete scheme). 
A $k$-point of $\Scr A$ is a formal finite sum 
\[ \sum_{v\in \abs C} \ch\lambda_v \cdot v \]
where $\ch\lambda_v$ is an element of the dual monoid $\mathfrak c_X$ and $\ch\lambda_v=0$ for all but finitely many $v$. 

\subsubsection{} 
The stratification described in \S \ref{sect:globstrat} takes here the following form:
For any $\mf P \in \Sym^\infty(\mathfrak c_X\sm 0)$ there is a natural map 
\[ \oo C^{\mf P} \into \Scr A, \] 
where the image consists of the $k$-points $\sum_{v\in \abs C} \ch\lambda_v\cdot v$ such that the unordered multiset of nonzero $\ch\lambda_v$, counted
with multiplicities, coincides with $\mf P$. 



\begin{prop}[{\cite[Proposition 3.5]{BNS}}]
\source{intersection-complexes-unramified-l-factors}
\notready 
\label{prop:Astrata}
\source{intersection-complexes-unramified-l-factors}
\ 
\begin{enumerate}
\item The maps $\oo C^{\mf P}\into \Scr A$ are locally closed embeddings,
and the collection of such embeddings over all $\mf P\in \Sym^\infty(\mathfrak c_X \sm 0)$ forms a stratification of $\Scr A$.
\item $\oo C^{\mf P'}$ lies in the closure of $\oo C^{\mf P}$
if and only if $\mf P$ refines $\mf P'$. 
\item The irreducible components of $\Scr A$ are in bijection with 
the closures of $\oo C^{\mf P}$ for $\mf P \in \Sym^\infty(\Prim(\mf c_X))$.
\end{enumerate}
\end{prop}

\begin{cor}[{\cite[Corollary 3.6]{BNS}}]
\source{intersection-complexes-unramified-l-factors}
\notready 
For $\ch\lambda \in \mathfrak c_X$, let 
\[ \Scr A^{\ch\lambda} \subset \Scr A \]
denote the subscheme whose $k$-points consist of all 
$\sum_{v\in \abs C} \ch\lambda_v \cdot v$
such that $\sum_v \ch\lambda_v = \ch\lambda$. 
Then $\Scr A^{\ch\lambda}$ is a connected component of $\Scr A$, and
this gives a bijection 
\[ \pi_0(\Scr A) \cong \mathfrak c_X. \] 
\end{cor}

\begin{rem}
\source{intersection-complexes-unramified-l-factors}
\notready\label{rem:Afree}
\source{intersection-complexes-unramified-l-factors}
If $\mathfrak c_X$ is a free monoid with basis $\{ \ch\nu_i \}$, 
then for $\ch\lambda = \sum_i N_i \ch\nu_i$ we have 
$\Scr A^{\ch\lambda} = \prod_i C^{(N_i)}$. The bases for the Zastava spaces
in \cite[\S 2.1]{BFGM} take this form.
\end{rem}



\subsection{Zastava model} \label{sect:localmodel}
\source{intersection-complexes-unramified-l-factors}

\emph{For the rest of this paper we assume that $B$ acts simply transitively on $X^\circ$.} 
This implies that $T_X = T$ and $\Lambda_X = \Lambda_G$, so we will use the notation
interchangeably. 

\medskip

We introduce a special case of the model used in \cite[Part III]{GN}, which is based 
on a general pattern pointed out by Drinfeld (see \cite[\S 4.2--4.4]{Drinfeld}). 
These are a generalization of the Zastava\footnote{``Zastava'' is Croatian for ``flag''.} spaces introduced by Finkelberg--Mirkovi\'c in \cite{FM, FFKM, BFGM}, and we will henceforth call them the Zastava model for $X$. 

\medskip

The Zastava model for $X$ is defined as 
\begin{equation}
    \sY = \Scr Y_{X} = \Maps_\gen(C, X/B \supset \pt),
\end{equation}
the stack of maps $C \times S \to X/B$ generically landing 
in $X^\circ/B=\pt$. 


Applying $\Maps(C,?/T)$ to the natural map $X/N \to X/\!\!/N $ induces a map 
\begin{equation}\label{mappi} \pi: \Scr Y \to \Scr A.
\source{intersection-complexes-unramified-l-factors}
\end{equation}
For $\ch\lambda \in \mf c_X$, let $\sY^{\ch\lambda}$ denote the preimage of 
$\sA^{\ch\lambda}$ under $\pi$. 

We show in Proposition~\ref{prop:YGr} below that $\sY$ is representable by a
scheme locally of finite type over $k$. This was predicted by 
Drinfeld \cite[Conjecture 4.2.3]{Drinfeld} in 
a more general setting. 

\begin{eg}
\source{intersection-complexes-unramified-l-factors}
\notready \label{eg:YHecke}
\source{intersection-complexes-unramified-l-factors}
Let $X = \mbb G_m \bs \GL_2$ where $\mbb G_m$ is embedded as $\left( \begin{smallmatrix} 1 & \\ & * \end{smallmatrix} \right)$. 
Then $\ch G_X = \ch G = \GL_2$ and $\ch\Lambda_X = \ch\Lambda_G = \mbb Z^2$ with 
standard basis $\ch\vareps_1,\ch\vareps_2$.
The $B$-orbits on $X$ are the same as $\mbb G_m$-orbits on $G/B=\mbb P^1$, so
there are three orbits: $\mbb G_m, \{0\}, \{\infty\}$. 
These correspond to $X^\circ$ and two colors $D^+, D^- \subset X$, respectively.
We have $\ch\nu_{D^+} = \ch\vareps_1,\, \ch\nu_{D^-} = -\ch\vareps_2$
and $\mf c_X = \mbb N^2$ is the free monoid generated by $\ch\nu_{D^+}, \ch\nu_{D^-}$. 
Note that $\ch\nu_{D^+} + \ch\nu_{D^-} = \ch\alpha$ is the simple coroot. 

Since $X$ is affine homogeneous, $\sM_X = \Bun_H = \Bun_1$ is the moduli stack
of line bundles. The Zastava model is $\sY=\Maps_\gen(C, \mbb G_m \bs \GL_2/B \supset \pt) 
= \Maps_\gen(C, \mbb G_m \bs \mbb P^1 \supset \pt)$. 
This is the stack parametrizing two line bundles $\Scr L, \Scr L'$ on $C$
and a fiberwise injective map of vector bundles $\sigma :\Scr L \into \Scr O_C \oplus \Scr L'$ 
such that \emph{both} coordinates $\sigma_1 : \Scr L \into \Scr O_C$ and $\sigma_2 : \Scr L \into \Scr L'$ are generically nonzero. 
Thus, $\sigma_1,\sigma_2$ are equivalent to two effective Cartier divisors 
$D_1, D_2$ which give $\Scr L = \Scr O(-D_1)$ and $\Scr L' = \Scr O(D_2-D_1)$. 
The condition that $\sigma=(\sigma_1,\sigma_2)$ is fiberwise injective is equivalent
to saying that the supports of $D_1,D_2$ are disjoint. 
Therefore, we have an identification
\[ \sY_{\mbb G_m \bs \GL_2} = \Sym C \oo\xt \Sym C = \bigsqcup_{(n_1,n_2)\in \mbb N^2} C^{(n_1)} \oo\xt C^{(n_2)}. \]
Meanwhile $\Scr A_{\mbb G_m \bs \GL_2} = \Sym C \xt \Sym C = \bigsqcup_{\mbb N^2} C^{(n_1)} \times C^{(n_2)}$. 
In this case $\sY \to \sA$ is the natural open embedding, and 
$\sY^{n_1\ch\nu_{D^+}+n_2\ch\nu_{D^-}} = C^{(n_1)} \oo\xt C^{(n_2)}$ is connected.
The map $\sY \to \sM_X$ forgetting the $B$-structure corresponds to 
$(D_1,D_2) \mapsto \Scr O(D_2-D_1)$. 

Note that for the embedding above of $\mbb G_m \into \GL_2$, 
the open $B$-orbit is the orbit of $\left( \begin{smallmatrix} 1 & 0 \\ 1 & 1 \end{smallmatrix}
\right) \in \mbb G_m \bs \GL_2$. In particular $X^\circ$ does not contain the identity coset. 
If we conjugate $\mbb G_m$ to the embedding $\left( \begin{smallmatrix} 1 & 0 \\ 1-a & a \end{smallmatrix}\right)$, then we may take the base point $x_0=1$. 
\end{eg}

\begin{eg}
\source{intersection-complexes-unramified-l-factors}
\notready \label{eg:YHeckePGL}
\source{intersection-complexes-unramified-l-factors}
In the example above we could instead replace $\GL_2$ by $G = \PGL_2$ and 
$H = \mbb G_m$ becomes the split torus inside $\PGL_2$. 
Then we still have $X = \mbb G_m\bs \PGL_2$ affine spherical with $\ch G_X = \ch G = \mathrm{SL}_2$ and two colors $D^+,D^-$. 
However now $\ch\Lambda_X = \ch\Lambda_G = \mbb Z \frac {\ch\alpha}2$ 
and $\ch\nu_{D^+} = \ch\nu_{D^-} = \frac {\ch\alpha}2$. 
The space $\Scr Y_{\mbb G_m \bs \PGL_2}$ still identifies with 
$\Sym C \oo\xt \Sym C$. However now $\Scr A_{\mbb G_m \bs \PGL_2} = \Sym C$
with $\Scr A^{n \frac{\ch\alpha}2} = C^{(n)}$. 
Then $\sY_{\mbb G_m \bs \PGL_2}^{n \frac{\ch\alpha}2} = \bigsqcup_{n_1+n_2=n} C^{(n_1)}\oo\xt C^{(n_2)}$ and the map $\sY \to \sA$ corresponds to addition of divisors.

Note that $\sM_X = \Bun_1$ and $C^{(n_1)}\oo\xt C^{(n_2)}$ maps to the component 
$\Bun_1^{n_2-n_1}$ of degree $n_2-n_1$ line bundles. 
Thus, while $\sY^{n \frac{\ch\alpha}2}$ is not connected, 
fixing a connected component of $\Bun_1$ determines the connected
component of $\sY^{n \frac{\ch\alpha}2}$.
\end{eg}


\subsection{Graded factorization property}
Note that our assumption on $X^\circ$ implies that $X^\circ/B=\pt$ is a dense open
substack of $X/B$. In the language of \cite[\S 4.2.1]{Drinfeld}, the stack $X/B$ is \emph{pointy}. Drinfeld observed that maps from a curve to a pointy stack will have local
behavior with respect to $C$ (compared to \cite{Drinfeld}, we are in the special
setting where we have a group $B$ acting on $X$, not just a groupoid).


Let $\ch\lambda = \ch\lambda_1 + \ch\lambda_2$ with $\ch\lambda_i \in \mathfrak c_X$ and let us denote by $\Scr A^{\ch\lambda_1} \oo\times \Scr A^{\ch\lambda_2}$ the open subset of the direct product $\Scr A^{\ch\lambda_1} \times \Scr A^{\ch\lambda_2}$ consisting of $\alpha_1,\alpha_2 : C \to (X/\!\!/N)/T$ such that
the supports of $C\sm \alpha_1^{-1}(\pt)$ and $C\sm \alpha_2^{-1}(\pt)$ are 
disjoint. 
We have a natural \'etale map $\Scr A^{\ch\lambda_1} \oo\times \Scr A^{\ch\lambda_2} \to \Scr A^{\ch\lambda}$.

\begin{prop}
\source{intersection-complexes-unramified-l-factors}
\notready \label{prop:factorization}
\source{intersection-complexes-unramified-l-factors}
The scheme $\Scr Y$ has the graded factorization property, in the sense
that there is a natural isomorphism 
\[ 
\Scr Y^{\ch\lambda}  \underset{\Scr A^{\ch\lambda}}\times 
(\Scr A^{\ch\lambda_1} \oo\xt \Scr A^{\ch\lambda_2})
\cong 
(\Scr Y^{\ch\lambda_1} \times \Scr Y^{\ch\lambda_2})|_{\Scr A^{\ch\lambda_1}\oo\times \Scr A^{\ch\lambda_2}}. 
\]
\end{prop}
\begin{proof}
Let $y_1, y_2 : C \times S \to X/B$ be $S$-points of $\Scr Y^{\ch\lambda_1}, \Scr Y^{\ch\lambda_2}$, respectively. 
Let $U_i = y_i^{-1}(\pt)\subset C\times S$; the condition that $(\pi(y_1),\pi(y_2))\in
\Scr A^{\ch\lambda_1} \oo\xt \Scr A^{\ch\lambda_2}$ is equivalent to 
requiring that $U_1 \cup U_2 = C\times S$. 
Then 
$y_1|_{U_1\cap U_2} \cong y_2|_{U_1\cap U_2} : U_1\cap U_2 \to \pt$
provides a gluing 
data for $y_1,y_2$ on the covering of $C\times S$ by $U_1$ and $U_2$.
Since $X/B$ is a stack, the gluing data descends to a map
$y : C\times S \to X/B$ that sends $U_1\cap U_2$ to $\pt$. 
This defines $y\in \Scr Y^{\ch\lambda}(S)$. 
The map in the opposite direction
is constructed in the same way and they are mutually inverse.
\end{proof}

We will henceforth use the notation 
$\Scr Y^{\ch\lambda_1} \oo\xt \Scr Y^{\ch\lambda_2}$ to denote
$(\Scr Y^{\ch\lambda_1} \times \Scr Y^{\ch\lambda_2})|_{\Scr A^{\ch\lambda_1}\oo\times \Scr A^{\ch\lambda_2}}$. 

\subsection{Global-to-Zastava yoga} 

We have a map $\Scr Y \to \Scr M_{X}$ by forgetting the $B$-structure. 
More precisely, we have an open embedding 
\begin{equation}
    \Scr Y \into \Scr M_{X} \underset{\Bun_G}\times \Bun_B.
\end{equation}

Although the natural map $\Bun_B \to \Bun_G$ is in general not smooth, 
it is smooth over a large enough open substack: 
consider $T$ as the Levi quotient of $B^-$. Let $\mf n^-$ denote the Lie algebra of $N^-$
viewed as a $T$-module. Define the open substack $\Bun^r_T \subset \Bun_T$ to consist
of those $T$-bundles $\Scr P_T$ for which 
$H^1(C, V \times^T \Scr P_T)=0$, for all $T$-modules $V$ which appear as
subquotients of $\mf n^-$. Let 
$\Bun_B^r$ be the preimage of $\Bun^r_T$ under the natural projection 
$\Bun_B \to \Bun_T$.

For $\ch\mu\in \ch\Lambda_T$, let $\Bun^{\ch\mu}_T$ denote the corresponding
connected component of $\Bun_T$ of degree $\ch\mu$, and 
let $\Bun^{\ch\mu}_B$ (resp.~$\Bun^{\ch\mu,r}_B$) be 
its preimage in $\Bun_B$ (resp.~$\Bun^r_B$). 
Note that by Riemann--Roch, $\Bun_B^{-\ch\mu,r} = \Bun_B^{-\ch\mu}$ 
if $\brac{\alpha_i,\ch\mu} > 2g-2$ for all simple roots $\alpha_i$, where $g$ is the genus of $C$. We say that $\ch\mu$ is ``large enough'' if $\brac{\alpha_i,\ch\mu}>N$
for all simple roots $\alpha_i$ and some $N \gg 0$.

\begin{lem}[{\cite{DS}, \cite[Lemma 3.7]{BFGM}, \cite[Lemma 14.2.1]{GN}}]
\source{intersection-complexes-unramified-l-factors}
\notready 
\label{lem:BunBtoGsm}
\source{intersection-complexes-unramified-l-factors}
The restriction of the map $\Bun_B \to \Bun_G$ to $\Bun^r_B$ is smooth. 
Any open substack $U\subset \Bun_G$ of finite type is contained in the image of 
$\Bun^{-\ch\mu}_B$ for $\ch\mu$ large enough, and the fibers of $\Bun^{-\ch\mu}_B\to \Bun_G$ over $U$ are geometrically connected. 
\end{lem}

Under our conventions, the composition $\Scr Y^{\ch\lambda} \to \Bun_B \to \Bun_T$ lands in the connected component $\Bun_T^{-\ch\lambda}$. 


\begin{cor}
\source{intersection-complexes-unramified-l-factors}
\notready\label{cor:Mcover}
\source{intersection-complexes-unramified-l-factors}
(i) The map $\Scr Y^{\ch\mu} \to \Scr M_{X}$
is smooth with geometrically connected fibers (when nonempty) for $\ch\mu$ large enough.

(ii)
Any $k$-point of $\Scr M_{X}$ lies in the image of $\Scr Y^{\ch\mu}$ for all $\ch\mu$ in a translate of $\ch\Lambda^\pos_G$. 
\end{cor}
\begin{proof}
We have an open embedding $\Scr Y^{\ch\mu} \into \Scr M_{X} \xt_{\Bun_G} \Bun_B^{-\ch\mu}$, so (i) follows from Lemma~\ref{lem:BunBtoGsm} 
by base change. To show (ii), we consider the fiber of 
$\Scr Y \to \Scr M_{X}$ on $k$-points:

A point in $\Scr M_X(k)$ is equivalent to a datum 
$(\Scr P_G, \sigma : C\to X\xt^G \Scr P_G)$. 
First we show that there exists some point in $\sY(k)$ that maps to $f$. 
By \cite{Steinberg}, there exists an open subscheme $U\subset C$
on which $\Scr P_G|_U$ can be trivialized. 
If we fix such a trivialization, then $\sigma|_U$ 
identifies with a section $U \to H\bs G$. Since $H\bs G$ is
spherical, $\sigma(U)$ intersects the open orbit of a 
Borel subgroup $g B g^{-1} \subset G$ for some $g\in G(k)$. 
Then $g$ defines a point in  
$(G/B)(k) \subset (G/B)(\Bbbk)$, where $\Bbbk=k(C)$. 
Using our fixed trivialization
of $\Scr P_G|_U$, we get a section 
\[ \Spec \Bbbk \overset{g}\to \Spec \Bbbk \times G/B \cong \Scr P_G^0|_{\Spec \Bbbk} \xt^G (G/B) 
\cong \sP_G|_{\Spec \Bbbk} \xt^G G/B, 
\]
which extends to a section $C \to \Scr P_G \xt^G G/B$ since 
$G/B$ is proper. The latter is equivalent to giving a 
$B$-structure $\Scr P_B$ on $\Scr P_G$. 
By construction $(\Scr P_B,\sigma)$ satisfies the generic condition 
for it to lie in $\Scr Y(k)$, and $(\Scr P_B,\sigma)$ maps to $(\sP_G,\sigma) \in \sM_X(k)$. 

Next let us fix an arbitrary lift of $(\Scr P_G,\sigma)$ to $(\Scr P_B^1,\sigma)\in \Scr Y(k)$.
Fix a trivialization of $\Scr P_B^1|_{\Spec \Bbbk}$, which also specifies a 
trivialization $\Scr P_G|_{\Spec \Bbbk} \cong \Scr P_G^0|_{\Spec \Bbbk}$.
With respect to this trivialization, we have a bijection of sets
\begin{equation}\label{e:liftBstructure}
\source{intersection-complexes-unramified-l-factors}
    H(\Bbbk)B(\Bbbk)/B(\Bbbk) \overset\sim\to \{\text{lift of } (\sP_G,\sigma) \text{ to a point in } \sY(k) \} 
\end{equation}
where the map is given by sending $h \in H(\Bbbk)B(\Bbbk)/B(\Bbbk)\subset (G/B)(\Bbbk)$ to 
the section
\[ \Spec \Bbbk \overset{h}\to \Spec \Bbbk \times G/B \cong \sP^0_G|_{\Spec \Bbbk} \xt^G(G/B)
\cong \sP_G|_{\Spec \Bbbk} \xt^G G/B \]
and uniquely extending to a section $C \to \sP_G \xt^G G/B$.
The point $B \in (G/B)(\Bbbk)$ is sent under \eqref{e:liftBstructure} to the lift $(\sP_B^1,\sigma)$.

We are concerned with the possible degrees of $\Scr P_B$ for 
lifts $(\Scr P_B,\sigma)\in \Scr Y(k)$ above $(\Scr P_G,\sigma)$.
For a simple root $\alpha$ of $G$, let $P_\alpha$ denote the corresponding minimal
parabolic in $G$. 
The image of $H\cap P_\alpha$ in $P_\alpha/\mf R(P_\alpha) = \PGL_2$ must contain 
a subgroup conjugate to $\left(\begin{smallmatrix} * & 0 \\ 0 & * \end{smallmatrix}\right)$
or $\left(\begin{smallmatrix} 1 & * \\ 0 & 1 \end{smallmatrix}\right)$.
Thus we can choose an identification $P_\alpha/B \cong \mbb P^1$ such that
the image of $(H\cap P_\alpha)(\Bbbk)$ in $\mbb P^1(\Bbbk)$ contains 
$\Bbbk^\times = k(C)^\times$. 
For $f \in (H\cap P_\alpha)(\Bbbk)$, let $\bar f$ denote its image in $(P_\alpha/B)(\Bbbk) \cong \mbb P^1(\Bbbk)$. 
For such a function $f$, let $(\sP_B,\sigma)$ denote the corresponding lift
of $(\sP_G,\sigma)$ under the bijection \eqref{e:liftBstructure}.
If $\bar f \in k(C)^\times$, then $\sP_B$ is of degree $-(\ch\mu_1 + N\ch\alpha)$, 
where $-\ch\mu_1$ is the degree of $\Scr P_B^1$ and 
$N$ is the degree of the divisor of zeros of $\bar f$.

By the Riemann--Roch theorem, we have a rational function $\bar f\in k(C)^\times$ 
with divisor of zeros of degree $N$ for any $N \gg 0$. 
Therefore for any $N\gg 0$, there is a lift $\Scr P_B$ of degree 
$-(\ch\mu_1+N\ch\alpha)$. 
We conclude that there exists a lift 
$(\Scr P_B,\sigma)\in \Scr Y^{\ch\mu}(k)$ for any $\ch\mu$ 
in a certain translate of $\ch\Lambda^\pos_G$.
\end{proof}


\subsubsection{} \label{sect:locglobyoga}
\source{intersection-complexes-unramified-l-factors}
It is not in general true that the natural map 
$\Scr Y^{\ch\lambda} \to \Scr M_{X}$ is smooth for 
arbitrary $\ch\lambda$. However, the following 
well-known argument (cf.~\cite[Theorem 16.2.1]{GN}) 
shows that any neighborhood of a point in $\Scr Y^{\ch\lambda}$
is smooth locally isomorphic to a neighborhood of a point in 
$\Scr M_{X}$:

Let $\ch\lambda\in \mathfrak c_X$ be arbitrary. 
Since $\ch\Lambda^\pos_G \subset \mf c_X$, we can always
find $\ch\mu \in \ch\Lambda^\pos_G$ large enough such that 
$\ch\lambda+\ch\mu$ is also large enough. 
Let $\Scr Y^{\ch\mu,0}$ denote the preimage of $\Scr M_X^0=\Bun_H$ under the smooth map $\Scr Y^{\ch\mu} \to \Scr M_{X}$.
Then $\Scr Y^{\ch\mu,0}$ is smooth and   
the first projection
\[  \Scr Y^{\ch\lambda} \times \Scr Y^{\ch\mu,0} \to \Scr Y^{\ch\lambda} \]
is smooth.
On the other hand, by the graded factorization property (Proposition~\ref{prop:factorization}), there is a natural \'etale map 
\[ \Scr Y^{\ch\lambda} \oo\xt \Scr Y^{\ch\mu,0} 
\to \Scr Y^{\ch\lambda+\ch\mu}. 
\]
We can compose this with the smooth map $\Scr Y^{\ch\lambda+\ch\mu}\to \Scr M_X$ to get a smooth map $\Scr Y^{\ch\lambda} \oo\xt \Scr Y^{\ch\mu,0} \to \Scr M_X$.
To summarize, we have constructed:


\begin{lem}
\source{intersection-complexes-unramified-l-factors}
\notready \label{lem:localglobalyoga}
\source{intersection-complexes-unramified-l-factors}
For any $\ch\lambda\in \mathfrak c_X$ and any $\ch\mu \in \ch\Lambda_G^\pos$ large enough, there is a correspondence 
\[ 
\Scr Y^{\ch\lambda} \leftarrow 
\Scr Y^{\ch\lambda} \oo\times \Scr Y^{\ch\mu,0} \rightarrow 
\Scr M_{X} 
\]
where the left arrow is smooth surjective, and the right arrow is smooth.
\end{lem}



\subsection{Stratification of the Zastava model}
We stratify $\sY^{\ch\lambda}$ according to the fine stratification of $\sM_X$:
for a partition $\ch\Theta \in \Sym^\infty(\mf c_X^- \sm 0)$
and $\ch\lambda\in \mf c_X$,
define the stratum 
\[ \Scr Y^{\ch\lambda,\ch\Theta} := \Scr Y^{\ch\lambda}
\underset{\Scr M_X}\times \Scr M_X^{\ch\Theta}. \]
Abbreviate $\sY^{\ch\lambda,\ch\theta} := \sY^{\ch\lambda,[\ch\theta]}$.

\begin{prop}
\source{intersection-complexes-unramified-l-factors}
\notready \label{prop:localstrata}
\source{intersection-complexes-unramified-l-factors}
The stratum $\Scr Y^{\ch\lambda,\ch\Theta}$ is a smooth locally closed
subscheme of $\Scr Y^{\ch\lambda}$.
\end{prop}
\begin{proof}
By Lemma~\ref{lem:localglobalyoga}, there exists $\ch\mu\in\ch\Lambda_X$
such that there is a smooth correspondence 
\[ \Scr Y^{\ch\lambda} \leftarrow \Scr Y^{\ch\lambda}\oo\times \Scr Y^{\ch\mu,0} \to \Scr M_X \]
where the left arrow is surjective.
Note that by definition of $\Scr M_X^{\ch\Theta}$, the preimage of 
$\Scr M_X^{\ch\Theta}$ in $\Scr Y^{\ch\lambda}\oo\times\Scr Y^{\ch\mu,0}$
is isomorphic to $\Scr Y^{\ch\lambda,\ch\Theta} \oo\times \Scr Y^{\ch\mu,0}$
since $\Scr Y^{\ch\mu,0}$ consists
of maps $C\to X^\bullet/B$ which can only define points in $\msf L^0 X/\msf L^+ G$ 
upon restriction to $\mf o_v$ for any $v\in \abs C$.
Therefore, we get a smooth correspondence
\begin{equation} \label{e:localglobalcorresp}
\source{intersection-complexes-unramified-l-factors}
    \Scr Y^{\ch\lambda,\ch\Theta} \leftarrow \Scr Y^{\ch\lambda,\ch\Theta} \oo\times \Scr Y^{\ch\mu,0} \to \Scr M_X^{\ch\Theta} 
\end{equation}
where the left arrow is still surjective. 
Now smoothness of $\Scr Y^{\ch\lambda,\ch\Theta}$ follows from smoothness
of $\Scr M_X^{\ch\Theta}$ (Lemma~\ref{lem:globstrata}).
\end{proof}

We call the collection of connected components of $\Scr Y^{\ch\lambda,\ch\Theta}$ the
\emph{fine stratification} of $\Scr Y^{\ch\lambda}$. By the smooth
correspondence \eqref{e:localglobalcorresp} above 
and Proposition~\ref{prop:globwhitney}, this is a Whitney stratification (in fact
the Zastava model is used in the proof of \emph{loc~cit.}).

Note that for a fixed $\ch\lambda$, many of the 
 $\Scr Y^{\ch\lambda,\ch\Theta}$ are empty.



\subsection{Relation to the affine Grassmannian} \label{sect:YGr} 
\source{intersection-complexes-unramified-l-factors}

Let $\Gr_{G,\Sym C} \to \Sym C$ denote the following version of the Beilinson--Drinfeld affine Grassmannian: an $S$-point consists of a relative effective Cartier divisor $D \subset C\times S$ and a $G$-bundle $\Scr P_G$ on $C\times S$ together with a trivialization
$\Scr P_G|_{C\times S \sm D} \cong \Scr P_G^0|_{C\times S\sm D}$ where 
$\Scr P_G^0$ is the trivial $G$-bundle. 
For any linear algebraic group $G$, the functor $\Gr_{G,\Sym C}$ is representable by an ind-scheme, ind-of finite type
over $\Sym C$ (cf.~\cite{BD}, \cite[Theorem 3.1.3]{Xinwen}).

In particular, we can consider the ind-scheme 
$\Gr_{B,\Sym C}$. Let 
$\Gr_{B,C^{(N)}}$ denote the preimage over $C^{(N)}$. 

\subsubsection{} \label{sss:YGr}
\source{intersection-complexes-unramified-l-factors}
Choose some $\delta \in \Lambda_X$ that lies on the interior of the
cone dual to $\Cal C_0(X)$, so $\brac{\delta,\ch\lambda}>0$
for any nonzero $\ch\lambda \in \Cal C_0(X)$.
Let $f_\delta \in k[X]^{(B)}$ denote the corresponding $\delta$-eigenfunction. 
Then $f_\delta$ induces a map $X/\!\!/N \to \mbb A^1$, which 
in turn induces a map $\Scr A \to \Sym C$ sending 
$\Scr A^{\ch\lambda} \to C^{(\brac{\delta,\ch\lambda})}$. 
We can map\footnote{
The effective Cartier divisor cut out by $f_\delta$ 
has the property that the complement of its support in $X$ equals $X^\circ$. 
In this guise, the map $\sY \to \Sym C$ we have constructed coincides with
the one described in \cite[Remark 4.2.6]{Drinfeld}.}

\begin{equation} \label{e:YGr}
\source{intersection-complexes-unramified-l-factors}
    \Scr Y^{\ch\lambda} \to \Gr_{B,C^{(\brac{\delta,\ch\lambda})}} 
\end{equation}
as follows:
let $y : C\times S \to X/B$ be an $S$-point of $\Scr Y$. 
Let $D \subset C \times S$ be the relative effective divisor corresponding
to the image of $\pi(y)\in \Scr A(S)$ under the map $\Scr A \to \Sym C$.
Since $\delta$ was chosen in the interior of the dual cone of $\Cal C_0(X)$, we have
$C\times S \sm D = y^{-1}(\pt) =: U$. 
Thus, $y$ defines a $B$-bundle $\sP_B$ on $C \xt S$ together with 
a section $U \to X^\circ \xt^B \sP_B \cong \sP_B$, i.e., a trivialization of 
$\sP_B|_U$. 
The datum $(D, \Scr P_B, \Scr P_B|_U \cong \Scr P_B^0|_U)$ defines an $S$-point 
of $\Gr_{B,\Sym C}$. 


\begin{prop}
\source{intersection-complexes-unramified-l-factors}
\notready \label{prop:YGr}
\source{intersection-complexes-unramified-l-factors}
The map \eqref{e:YGr} is a closed embedding. 
Moreover, the stack $\Scr Y$ is representable by a scheme locally of finite type over $k$, and for fixed $\ch\lambda\in \mathfrak c_X$, the scheme 
$\Scr Y^{\ch\lambda}$ is of finite type over $k$. 
Consequently, the map $\pi : \Scr Y \to \Scr A$ is schematic of finite type. 
\end{prop}
\begin{proof}
First we show that \eqref{e:YGr} is a closed embedding.
Fix a map $S \to \Gr_{B,\Sym C}$ corresponding to a pair 
$(D,\Scr P_B)$ and a trivialization of $\Scr P_B$ away from $D$. 
A trivialization of $\Scr P_B$ on $C\times S\sm D$ is equivalent to a section 
$\sigma_0:C \times S\sm D \to \Scr P_B \cong X^\circ \xt^B \sP_B$. 
The fiber of $\Scr Y \to \Gr_{B,\Sym C}$ over $S$
 parametrizes commutative diagrams 
\[ 
\begin{tikzcd}
C \times S' \sm D' \ar[d,hook] \ar[r, "{\sigma_0}"] & X^\circ \overset B\times\Scr P_B \ar[d, hook]  \\ 
C \times S' \ar[r, "\sigma"] & X \overset B\times \Scr P_B  
\end{tikzcd}
\]
where $S'$ is an $S$-scheme, $D':= D\times_S S'$, and $\sigma$ is a 
section over $C\times S$. 
Observe that $\sigma$ is uniquely determined by $\sigma_0$. 
By Lemma~\ref{lem:extendsection} below, the condition that 
$\sigma_0$ extends to $\sigma$ is closed in $S$. 

We have shown that $\Scr Y$ is representable by an ind-scheme
ind-closed in $\Gr_{B,\Sym C}$. On the other hand, we have a map
$\Scr Y \to \Bun_B$ whose fiber over an $S$-point $\Scr P_B \in \Bun_B(S)$
is open in the space of sections $C\times S \to X \times^B \Scr P_B$ over $C\times S$.
This space is representable by a scheme locally of finite type over $k$ (\cite[Theorem 5.23]{FGA-explained}). Since $\Bun_B$ is an algebraic stack, we conclude that 
$\Scr Y$ is an algebraic stack representable by an ind-scheme. Hence 
$\Scr Y$ is representable by a scheme. 
For fixed $\ch\lambda \in \mathfrak c_X$, we now know that $\Scr Y^{\ch\lambda}$ is a closed subscheme of $\Gr_{B,C^{(\brac{\delta,\ch\lambda})}}$, which is ind-of finite type. It follows that $\Scr Y^{\ch\lambda}$ is of finite type over $k$. The other assertions all follow.  
\end{proof}

\begin{lem}
\source{intersection-complexes-unramified-l-factors}
\notready \label{lem:extendsection} 
\source{intersection-complexes-unramified-l-factors}
Let $S$ be a test scheme and $D \subset C\xt S$ a relative effective divisor.
Let $\Scr X$ be a scheme affine of finite presentation over $C\xt S$.
Suppose that there exists a section 
$\sigma : C\xt S \sm D \to \Scr X$ over $C\xt S$. 
Then the functor sending $S'$ to the set of 
maps $S'\to S$ such that $\sigma$ extends to a regular map on $C\xt S'$
is representable by a closed subscheme of $S$. 
\end{lem}
\begin{proof} The map $\sigma$ is equivalent to a map of 
$\Scr O_{C\xt S}$-algebras
$\Scr O_{\Scr X} \to \Scr O_{C\xt S\sm D}$. 
Given $S' \to S$, the condition that 
$\sigma$ extends to $C\xt S'$ is equivalent to 
requiring the image of 
$\Scr O_{\Scr X} \to \Scr O_{C\times S\sm D}$ to land in $\Scr O_{C\times S'} \subset \Scr O_{C\times S' \sm D'}$
after base change to $S'$, where $D':= D\xt_S S'$. 
The claim is local in $S$, so we may assume that
$\Scr O_{\Scr X}$ is surjected onto by 
$\Sym_{\Scr O_{C\xt S}}(\Scr E^\vee)$ for some
vector bundle $\Scr E$ on $C\xt S$. 
Then we just need the composed $\Scr O_{C\times S}$-linear map $\Scr E^\vee \to \Scr O_{C\times S \sm D} \to \Scr O_{C\times S\sm D}/\Scr O_{C\times S}$ to vanish after base change to $S'$.
Since $\Scr E^\vee$ is coherent, the image of this map 
is contained in a submodule $\Scr F = \Scr O(m\cdot D)/\Scr O_{C\times S}$
for some integer $m \ge 0$.  
The projections $p: C\times S \to S,\, p':C\times S'\to S'$ are proper, and we are considering
when an element in $H^0 p'_*( \Scr E \otimes_{\Scr O_{C\times S}} \Scr F \otimes_{\Scr O_S} \Scr O_{S'})$ vanishes. 
Note that $p_*(\Scr E \otimes_{\Scr O_{C\times S}} \Scr F)$ is finite locally free
as an $\Scr O_S$-module.
By cohomology and base change 
we are reduced to asking when an element of $p_*(\Scr E \otimes_{\Scr O} \Scr F)$
vanishes after base change to $S'$. 
This is a closed condition on $S$.
\end{proof}

\begin{rem}[Open curves]
\source{intersection-complexes-unramified-l-factors}
\notready \label{rem:opencurve}
\source{intersection-complexes-unramified-l-factors}
The graded factorization property of $\Scr Y^{\ch\lambda}$ implies that 
the geometry of $\Scr Y^{\ch\lambda}$ is purely local
with respect to the curve $C$. Therefore, we could define 
\[ \Scr Y(C) = \Maps_\gen(C, X/B \supset \pt) \]
for any smooth curve $C$ (not necessarily projective), and all 
the same properties would still hold.
For example, in \cite{FM}, \cite{Drinfeld}, the affine curve $C = \mbb A^1$ is used.
\end{rem}

\subsubsection{Beauville--Laszlo's theorem} \label{sect:BLthm}
\source{intersection-complexes-unramified-l-factors}
Let $S$ be an affine scheme and $D$ a closed affine subscheme of $C\times S$.
Denote by $\wh C_D$ the formal completion of $C\times S$ along $D$ and by
$\wh C'_D$ the spectrum of the ring of regular functions on $\wh C_D$ (so $\wh C_D$ is an ind-affine formal scheme and $\wh C'_D$ is the corresponding true scheme). 
Let $\wh C^\circ_D := \wh C'_D \sm D$ denote the open subscheme.

There are maps $\hat p: \wh C_D \to C\times S$ and $i: \wh C_D \to \wh C'_D$. 
We will implicitly use the following fact in what follows:

\begin{prop}[{\cite[Proposition 2.12.6]{BD}}]
\source{intersection-complexes-unramified-l-factors}
\notready \label{prop:formalext}
\source{intersection-complexes-unramified-l-factors}
There exists a unique map 
$p : \wh C'_D \to C\times S$ such that $\hat p = p\circ i$. 
\end{prop}


To justify that $\Scr Y$ is of local nature, we record 
the following consequence of the globalized version of Beauville--Laszlo's theorem (cf.~\cite[Theorem 2.12.1]{BD}, \cite{BLa}). 
Let $S$ be an affine scheme and $D\subset C\times S$ a relative effective divisor.  
Proposition~\ref{prop:formalext} implies that there exists a map $p:\wh C'_D\to C\times S$. 

\begin{lem}
\source{intersection-complexes-unramified-l-factors}
\notready \label{lem:B-L}
\source{intersection-complexes-unramified-l-factors}
Let $X$ be any affine scheme with an action of an algebraic group $B$
such that $X/B$ is pointy, i.e., $X$ has an open $B$-orbit $X^\circ$ with $X^\circ/B=\pt$.
Let $C,S$ and $D$ be as above. 

Then there is a natural equivalence between the following categories:
\begin{enumerate}
\item the groupoid of $(\Scr P_B, \sigma)$ where $\Scr P_B$ is a $B$-bundle on $C\times S$ and a section $\sigma : C\times S \to X\times^B \Scr P_B$ 
that sends $C\times S\sm D$ to $X^\circ \times^B \Scr P_B$,
\item the groupoid of $(\hat{\Scr P}_B, \hat \sigma)$ where $\hat{\Scr P}_B$ is a $B$-bundle on $\wh C'_D$ and a section $\hat\sigma: \wh C'_D \to X\times^B \hat{\Scr P}_B$ that sends $\wh C^\circ_D$ to $X^\circ \times^B \hat{\Scr P}_B$.
\end{enumerate}
\end{lem}

\begin{proof}
The functor from (i) to (ii) is just pullback along $p$. 
To define the functor from (ii) to (i) we descend along the covering 
$\wh C'_D \sqcup (C\times S \sm D) \to C\times S$. The justification
for this is Beauville--Laszlo's theorem (cf.~\cite[Theorem 2.12.1]{BD}). 
First, $\hat \sigma$ induces 
a section $\wh C^\circ_D \to X^\circ \times^B \hat{\Scr P}_B \cong \hat{\Scr P}_B$.
Then we can ``descend'' $\hat{\Scr P}_B$ to a $B$-bundle $\Scr P_B$ on $C\times S$ with a section $C\times S\sm D \to \Scr P_B$ (which
is equivalent to a trivialization of $\Scr P_B|_{C\times S\sm D}$).
The section $\hat\sigma$ is equivalent to a map of quasicoherent
$\Scr O_{\wh C'_D}$-algebras $\Scr O_{X\times^B \hat{\Scr P}_B} \to \Scr O_{\wh C'_D}$ since $X$ is affine. 
Again by \cite[Theorem 2.12.1]{BD}, this descends 
to a map of quasicoherent $\Scr O_{C\times S}$-algebras $\Scr O_{X\times^B \Scr P_B} \to \Scr O_{C\times S}$ such that the restriction to 
$C\times S\sm D$ factors through the trivialization 
$\Scr O_{X^\circ \times^B \Scr P_B}\cong \Scr O_{\Scr P_B} \to \Scr O_{C\times S\sm D}$.
By construction, the two functors are mutually inverse.
\end{proof}


\subsection{Theorem of Grinberg--Kazhdan, Drinfeld} 

We now justify why $\sM_X$ and $\sY$ are indeed ``models'' for the formal arc space
$\msf L^+ X$. Since $\sM_X$ and $\sY$ are smooth-locally isomorphic,
it suffices to explain the latter. 

\begin{defn}
\source{intersection-complexes-unramified-l-factors}
\notready\label{def:formalmodel} 
\source{intersection-complexes-unramified-l-factors}
A finite type formal model of $\msf L^+ X$ at $\gamma_0\in \msf L^+ X(k)$ is the formal completion $\wh{Y}_y$ of a $k$-scheme of finite type $Y$ at a point $y\in Y$ equipped with an isomorphism of formal schemes
\begin{equation} \label{finite model}
\source{intersection-complexes-unramified-l-factors}
\wh{\msf L^+ X}_{\gamma_0} \simeq \wh{Y}_y \times \wh{\mbb A}^\infty,
\end{equation}
where $\wh{\mbb A}^\infty$ is the product of countably many copies of the formal disk $\on{Spf} k\tbrac t$.
\end{defn}


Since $X/B \supset \pt$ is a pointy stack, 
Drinfeld's proof \cite[\S 4.2-4.3]{Drinfeld} of the Grinberg--Kazhdan theorem 
essentially shows that the scheme $\sY$ (which we have shown is a disjoint union
of finite type schemes) explicitly satisfies the following:

\begin{thm}[Grinberg--Kazhdan, Drinfeld]
\source{intersection-complexes-unramified-l-factors}
\notready  \label{thm:DGK}
\source{intersection-complexes-unramified-l-factors}
Fix an arc $\gamma_0 : \Spec k\tbrac t \to X$ in $\msf L^+ X(k)$
such that $\gamma_0(\Spec k\lbrac t) \subset X^\circ$. 
Then there exists a point $y \in \Scr Y(k)$ 
such that the formal completion of $\Scr Y$ at $y$ is a finite type formal model of $\msf L^+ X$ at $\gamma_0$. 

More precisely, if $\gamma_0$ belongs to the stratum $\msf L^{\ch\theta} X$, for $\ch\theta\in \mf c_X^-$, we can take $y$ to be the point $t^{\ch\theta}$ in the central fiber $\msf Y^{\ch\theta,\ch\theta}$ over any point $v\in \abs C$.
\end{thm}


The central fiber $\msf Y^{\ch\theta}$ 
over a point $v\in \abs C$ is the fiber of $\sY^{\ch\theta}\to \sA^{\ch\theta}$ over the 
``diagonal divisor'' $\ch\theta \cdot v \in \sA^{\ch\theta}(k)$.
We define $\msf Y^{\ch\theta,\ch\theta} := \sY^{\ch\theta} \xt_{\sM_X} \sM_X^{\ch\theta}$. 
Then $\msf Y^{\ch\theta}$ is naturally a subvariety of $\Gr_B$ (see \S\ref{def:centralfiber}), 
and $t^{\ch\theta}$ denotes the corresponding point in $\Gr_B$. The significance of
this point will become evident in Corollary~\ref{cor:Ytheta}(iii).

\smallskip

The statements all follow from the proof of \cite[\S 4]{Drinfeld}. 
We also give the same argument, with some notational changes, in the proof of Theorem~\ref{thm:DGKfamily}.


\begin{rem}
\source{intersection-complexes-unramified-l-factors}
\notready
The point $y : C \to X/B$ can be chosen so that $y^{-1}(\pt) = C\sm v$ for a single
point $v\in \abs C$. However it is essential, for the theorem to hold, that $\sY$ contains
maps with multiple points of $C$ mapping to $(X\sm X^\circ)/B$. 
\end{rem}
